\documentclass[../../../../master.tex]{subfiles}
\begin{document}

\monographleaf{Neural Architectures, Representations \& Regularization}{Computer Vision, Image Formation \& 3D Scene Understanding}{Image Formation, Fourier Filtering, Pyramids \& 3D Geometry}{Image Formation, Lenses, Color Spaces \& Fourier Image Filtering}{\bookcite{foundations_computer_vision_mit}{Ch. 5--8, 15--18, 20}, \bookcite{dive_into_deep_learning_d2l}{Ch. 7}}

\section{Problem Formulation \& Unified Notation}

Computer vision inverts the physical \emph{forward image formation} process: a 3D scene illuminated by light sources reflects radiance onto a sensor through optics, producing a discrete 2D array $I[n_y, n_x, c]$ \bookcite{foundations_computer_vision_mit}{Ch. 5, pp. 99--113}. Classical and neural vision pipelines analyze and transform this signal via linear shift-invariant (LSI) filtering, spectral frequency decomposition, and sampling theory \bookcite{foundations_computer_vision_mit}{Ch. 15--16, pp. 267--328}, \bookcite{dive_into_deep_learning_d2l}{Ch. 7, pp. 273--285}.

\begin{crossbook}[Cross-Book Notation: Optics, Convolution vs. Cross-Correlation]
\begin{center}
\footnotesize
\begin{tabularx}{\linewidth}{@{} >{\raggedright\arraybackslash}p{2.5cm} >{\raggedright\arraybackslash}X >{\raggedright\arraybackslash}X >{\raggedright\arraybackslash}X @{}}
\toprule
\textbf{Concept} & \textbf{MIT Vision Book} & \textbf{D2L (Zhang et al.)} & \textbf{Unified Monograph} \\
\midrule
2D Image & $f[n,m]$ or $\ell(x,y)$ & $\mat{X} \in \R^{n_h \times n_w}$ & $I[y,x]$ or $I(\vecb{x})$ \\
Linear Filtering & Convolution $f \circ h$ vs.\ correlation & Cross-correlation $\mat{X} \star \mat{K}$ & $(I * K)[y,x]$ (flipped), $(I \star K)$ (corr.) \\
2D DFT Frequency & $(u, v)$ or $(\omega_x, \omega_y)$ & Spatial kernel size $k_h \times k_w$ & $\hat{I}(u,v) = \mathcal{F}\{I\}(u,v)$ \\
\bottomrule
\end{tabularx}
\captionof{table}{Notation reconciliation across \bookcite{foundations_computer_vision_mit}{pp. 9--15, 267--295} and \bookcite{dive_into_deep_learning_d2l}{Ch. 7, pp. 273--280}.}
\end{center}
\end{crossbook}

\section{Core Mathematical Derivations}

\subsection{Pinhole Optics, Thin-Lens Equation \& Color Spaces}

\begin{coremath}[Perspective Projection, Thin-Lens Law \& Radiometry]
Under an ideal pinhole at distance $f$ (focal length) from the image plane, a 3D point $(X, Y, Z)\T$ projects to $(x, y)\T$ by similar triangles \bookcite{foundations_computer_vision_mit}{Ch. 5, pp. 103--108}:
\begin{equation}
  x = f \frac{X}{Z}, \qquad y = f \frac{Y}{Z}.
\end{equation}
To gather finite photon flux without blurring, a thin convex lens of focal length $f$ focuses rays from object depth $z_o$ onto sensor distance $z_i$ via the \emph{Gaussian thin-lens equation} \bookcite{foundations_computer_vision_mit}{Ch. 6, pp. 114--126}:
\begin{equation}
  \frac{1}{z_o} + \frac{1}{z_i} = \frac{1}{f}.
\end{equation}
Out-of-focus points produce a circle of confusion of diameter $c \propto \frac{A}{z_o}|z_o - z_{\mathrm{focus}}|$ where $A = f / N$ is aperture diameter ($N$: f-number). For a Lambertian surface with albedo $\rho(\vecb{x})$ and normal $\vecb{n}$ under directional light $\vecb{l}$, reflected radiance is $L_o = \rho(\vecb{x})\max(0, \vecb{n}\T\vecb{l}) E_0$ \bookcite{foundations_computer_vision_mit}{Ch. 5, pp. 100--102}. Trichromatic sensors integrate spectral power $S(\lambda)$ against cone/filter curves $c_k(\lambda)$, giving linear camera responses $r_k = \int S(\lambda) c_k(\lambda)\,\mathrm{d}\lambda$ before nonlinear gamma encoding ($I_{\mathrm{sRGB}} \approx I_{\mathrm{lin}}^{1/2.2}$) \bookcite{foundations_computer_vision_mit}{Ch. 7--8, pp. 147--178}.
\end{coremath}

\subsection{Linear Shift-Invariant Filtering, 2D Fourier Analysis \& Nyquist Sampling}

\begin{coremath}[2D Convolution Theorem, Derivative Filters \& Nyquist-Shannon Limit]
Any Linear Shift-Invariant (LSI) operator $\mathcal{L}$ on images is uniquely determined by its impulse response $h$: $(I * h)[y,x] = \sum_{i,j} I[y-i, x-j]\, h[i,j]$ \bookcite{foundations_computer_vision_mit}{Ch. 15, pp. 267--285}. Because complex exponentials $e^{i 2\pi(ux+vy)}$ are eigenfunctions of all LSI operators, the 2D Discrete Fourier Transform diagonalizes convolution into pointwise multiplication \bookcite{foundations_computer_vision_mit}{Ch. 16, pp. 292--325}:
\begin{equation}
  \hat{I}[u,v] = \sum_{y=0}^{H-1}\sum_{x=0}^{W-1} I[y,x]\,e^{-i 2\pi\left(\frac{ux}{W} + \frac{vy}{H}\right)}, \qquad \mathcal{F}\{I * h\} = \hat{I} \odot \hat{h}.
\end{equation}
Gaussian smoothing $G_\sigma(x,y) = \frac{1}{2\pi\sigma^2}e^{-(x^2+y^2)/(2\sigma^2)}$ has a Gaussian spectrum $\hat{G}_\sigma(u,v) = e^{-2\pi^2\sigma^2(u^2+v^2)}$, acting as an isotropic low-pass filter \bookcite{foundations_computer_vision_mit}{Ch. 17, pp. 331--348}. Image gradients $\nabla I = (\partial_x I, \partial_y I)\T$ and edges (zero-crossings of Laplacian $\Delta I = \partial_{xx} I + \partial_{yy} I$) are regularized via Derivative-of-Gaussian ($\nabla G_\sigma * I$) and Sobel kernels \bookcite{foundations_computer_vision_mit}{Ch. 18, pp. 352--380}. By the \emph{Nyquist--Shannon sampling theorem}, sampling with spacing $\Delta x$ aliases frequencies $|u| > \frac{1}{2\Delta x}$; anti-aliasing low-pass prefiltering is therefore mandatory before downsampling \bookcite{foundations_computer_vision_mit}{Ch. 20, pp. 409--428}.
\end{coremath}

\begin{intuition}[Why CNNs Compute Cross-Correlation and Why Depthwise Separability Works]
Deep learning libraries implement $(I \star K)$ without flipping the kernel because learnable weights absorb the flip \bookcite{dive_into_deep_learning_d2l}{Ch. 7, pp. 274--276}. Moreover, a 2D Gaussian filter (and many edge kernels) is rank-1 separable: $G_\sigma(x,y) = g_\sigma(x)g_\sigma(y)$, reducing $O(k^2 HW)$ filtering cost to $O(2k HW)$ \bookcite{foundations_computer_vision_mit}{Ch. 15, pp. 282--285}.
\end{intuition}

\section{Visual Diagnostics \& Geometry}

\begin{figure}[htbp]
\centering
\begin{tikzpicture}[scale=0.85, every node/.style={font=\scriptsize}]
  % Pinhole diagram (left)
  \draw[thick, monographNavy] (0,-1.4) -- (0,-0.25) (0,0.25) -- (0,1.4);
  \node[below] at (0,-1.45) {Pinhole ($Z=0$)};
  \draw[dashed, gray] (-3.2,0) -- (2.4,0);
  \draw[very thick, monographTeal, ->] (-2.6,0) -- (-2.6,1.1) node[above] {Scene $(X,Y,Z)$};
  \draw[very thick, monographCrimson, ->] (1.6,0) -- (1.6,-0.68) node[below] {Sensor $(x,y)$};
  \draw[monographIndigo, thick] (-2.6,1.1) -- (1.6,-0.68);
  \draw[thick] (1.6,-1.2) -- (1.6,1.2) node[above] {Image plane ($z=f$)};

  % Fourier spectrum (right)
  \begin{scope}[xshift=6.5cm]
    \draw[->] (-2.2,0) -- (2.3,0) node[right] {$u$ (freq)};
    \draw[->] (0,-0.2) -- (0,2.1) node[above] {$|\hat{G}_\sigma(u)|$};
    \draw[monographNavy, very thick, domain=-2.0:2.0, samples=60] plot (\x, {1.8*exp(-2.2*\x*\x)});
    \draw[monographCrimson, dashed, thick] (-1.1,0) -- (-1.1,1.6) (1.1,0) -- (1.1,1.6);
    \node[monographCrimson, above] at (1.1,1.6) {Nyquist $\pm\frac{1}{2\Delta x}$};
  \end{scope}
\end{tikzpicture}
\caption{Left: perspective pinhole geometry $y = fY/Z$. Right: low-pass Gaussian spectrum attenuating frequencies above the Nyquist limit to prevent aliasing.}
\end{figure}

\begin{pitfall}[Gamma Nonlinearity \& Unfiltered Strided Downsampling]
Applying linear blur, interpolation, or physically-based shading in gamma-compressed sRGB space introduces dark halos and hue shifts; always linearize ($I_{\mathrm{lin}} = I_{\mathrm{sRGB}}^{2.2}$) before physical operations \bookcite{foundations_computer_vision_mit}{Ch. 7, pp. 152--156}. Similarly, strided pooling or `stride=2` convolutions without anti-aliasing low-pass filters fold high-frequency edges into spurious low-frequency Moir\'e patterns, breaking shift invariance \bookcite{foundations_computer_vision_mit}{Ch. 20, pp. 412--425}.
\end{pitfall}

\end{document}
